\section{讨论与局限}
\label{sec:discussion}

\subsection{选择器层误修复对无人值守 CI 自愈意味着什么}
LLM 自愈最常见的部署模式是高度自主化的：一个失败的端到端测试触发 LLM 给出补丁，CI 自动提交补丁或将其作为 PR 评论发出。由 F1 的对抗性上界粗略换算，在不带行为验证器（oracle）的此类部署中，约 15\%（粗略上界估计，非实测部署率）的真实 Playwright 断裂存在被朴素修复器静默改写到错误元素的风险。其危险并不在于“修复器让构建失败”——构建失败本就是测试应承担的职责——而在于“修复器使构建通过，但通过对的是错误的目标”。需再次强调，F1 衡量的是选择器解析层面的错指风险；它不是端到端运行时通过率。

\subsection{为什么软提示不是可审计安全边界}
Qwen 的 \textsc{Vanilla} 与 \textsc{Abstain} 两种提示之间逐例完全等同（表~\ref{tab:falseheal-paired}），说明小型本地代码模型并未把“必要时放弃”的软指令转化为可观察行为。GPT-4o 的结果更复杂：放弃提示将误修复从 63/75 降至 51/75，但残留误修复仍占 68.0\%。一种可能的解释是，软指令需与“产出像样答案”这一更强先验竞争，且模型越强越可能部分遵守该指令（此为机理层面的推测，本文未做训练动力学层面的验证）。对工具设计的启示仍然清晰：反误修复的安全保障必须可审计地置于模型\emph{之外}，例如确定性的检索分数阈值、锚点类别归属规则，以及一个能够证伪通过补丁的行为验证器。

\subsection{确定性杠杆的有限含义}
F2 不构成统计显著的性能 claim。它的价值在于给出一条低成本、零回归的工程线索：当系统知道某类锚点更稳定时，把约束放进候选检索与输出后置过滤，比只要求模型“优先考虑稳定锚点”更可审计。本文据此把 F2 作为辅助性经验，而非主要科学贡献。

\subsection{局限}
\textit{模型范围。}主实验使用 \texttt{qwen2.5-coder:7b}，修订版补充了 \texttt{gpt-4o} 在相同 F1 协议上的对照。该对照用于检验“软提示是否因小模型指令遵循不足而失效”：结果显示 GPT-4o 确能把误修复从 84.0\% 降至 68.0\%，但仍不能把风险降到可无人值守接受的水平。该结论仍只覆盖一个前沿模型、一个数据集与一次贪心解码设置；因此本文不声称已穷尽模型规模、训练方式或供应商差异。
\textit{基准数据集范围。}主结果来自单一 React+Playwright 代码库 koenig；F3 中已扩展至 payload，但其可达率受到 BEM 类名盲区的限制。
\textit{静态修复代理与运行时通过率。}77.3\% / 62.4\% 这两个 L3 数字反映的是选择器解析器层面的精确一致性。其与真实 Playwright 运行时通过率之间的换算关系，需要在 ReproBreak 提供的逐 commit 复现环境中通过 Docker 重放打过补丁的测试加以建立。\textit{运行时 pilot 仅用于证明鸿沟存在，而非估计通过率。}作为最小可行性检查，我们对 3 个 v1 静态精确匹配用例（ids 615, 702, 703）做了 Docker 运行时重放：三者的 fixed baseline 均通过；替换为 HealReact 提议选择器后，2/3 运行通过，1/3 运行失败（id 702，等待 \texttt{getByTestId('color-picker-toggle')} 超时）。这一小规模 pilot 不能估计全量通过率，但说明静态精确匹配既有运行时可迁移的案例，也会因可见性/时序/上下文差异而失败；因此本文继续把 77.3\% / 62.4\% 标为静态代理，而不宣称端到端通过率。

\subsection{有效性威胁}
78.7\% 这一数字依赖解析器语义。\texttt{resolve\_locators.py} 把动态 JS 表达式属性值视为可能匹配、对复合 CSS 使用祖先过近似；这两项决策在 triage 场景下合理，但在边界情况下偏松。我们在解析器假设审计中撤回了早期“可达率上限 98\%”的虚抬声明，在收紧解析器假设后以 80.6\% 作为保守数字。由于解析器是本文所有“正确匹配/误修复/可达率”判定的事实 oracle，其判定边界会直接影响这些数字，我们做了一次 strict/permissive 双假设的敏感性扫描（表~\ref{tab:resolver-sens}）。本文头条的 80.6\%（75/93）来自 permissive 假设；若关闭上述两处过近似改用 strict 假设，可达率降至 54.8\%（51/93）。因此 L1 静态可达率的真实区间应理解为 [54.8\%, 80.6\%]，80.6\% 是上界。此外，permissive 下 75 个可达用例中有 46 个为多匹配（解析器按 \texttt{(componentFile, line)} 顺序取首个命中消歧），未解析用例从 permissive 的 18 个增至 strict 的 42 个。对抽样用例用真实 DOM 快照交叉校验解析判定，仍留作后续工作。

\begin{table}[t]
\centering
\caption{解析器 strict/permissive 敏感性（koenig，$n=93$）。permissive 为本文所采用的假设；strict 关闭“动态属性值视为可能匹配”与“同文件祖先近似”两处过近似。}
\label{tab:resolver-sens}
\small
\resizebox{\columnwidth}{!}{%
\begin{tabular}{lrr}
\toprule
指标 & permissive（本文） & strict（下界） \\
\midrule
new\_locator 可达 & 75 / 93 (80.6\%) & 51 / 93 (54.8\%) \\
多匹配（$\ge2$ 命中） & 46 & 0 \\
未解析（0 命中） & 18 & 42 \\
\bottomrule
\end{tabular}}
\end{table}

\subsection{未来工作}
本文的诊断结论自然延伸出两类可推进的工程化方向：其一，将静态基底向跨锚点文化（BEM 模板字面量、CSS Modules 等）扩展，使可达率在不同 React 项目中保持稳定；其二，沿 \S\ref{sec:findings} 给出的经验性必要性，构建并集成一个基于网络行为重放的轻量级验证器，作为 LLM 修复结果的下游守门。
